% SPDX-License-Identifier: LPPL-1.3c
% Copyright (C) 2026 Christophe Jorssen
% Author and Current Maintainer: Christophe Jorssen <christophe.jorssen@gmail.com>
% LPPL maintenance status: maintained
% This file is part of luafemm. See LICENSE and LICENSES.md for its terms.
\input tikz.tex
\input pgfplots.tex
\usetikzlibrary{femm}
\usepgfplotslibrary{femm}
\pgfplotsset{compat=1.18}
\tikzset{confined iron/.style={draw,even odd rule,
  femm/region={material=pure-iron,ideal domain=true}}}
% Excitation precedes the problem key: it must be queued for this picture.
\tikzpicture[femm/excitation={ampere turns=200},
  femm/problem={field model=ideal,depth=10,mesh size=3},scale=.6]
  \path[confined iron] (-20,-20) rectangle (20,20)
    (-10,-10) rectangle (10,10);
  \path[femm/profile={name=ideal-section}] (10,0)--(20,0);
  \pic[femm/lines=7] {femm field};
  \directlua{assert(luafemm.current.ideal_sources[1][3]==200)}
\endtikzpicture
\par
\femmflux[profile=ideal-section]
\par
\tikzpicture[femm/problem={field model=ideal,depth=10,mesh size=4},scale=.65]
  \node[femm/tapered toroid={ampere turns=0,magnet span=30,
    magnet material=ndfeb-40-mgoe},rotate=13,anchor=core center] (Ring) at (3,2) {};
  \path[femm/profile={name=pole-section}]
    (Ring.upper pole inner)--(Ring.upper pole outer);
  \pic[femm/lines=9] {femm field};
  \pic[femm/mean={component=Ring,kind=flux,profile=ring,samples=29}]
    {femm mean path};
  \directlua{
    local p=require('luafemm-profiles').sample('ring')
    assert(rawlen(p)==29 and p[1].B>0)
    local m=luafemm.current
    assert(rawlen(m.ideal_loops)==2)
  }
\endtikzpicture
\par
\femmflux[profile=pole-section]
\par
% The first profile still refers to the preceding model, after a new solve.
\femmflux[profile=ideal-section]
\par
\tikzpicture
\axis[width=7cm,height=4cm]
\femmplot[profile=ring,component=circulation]
\endaxis
\endtikzpicture
\bye
